% Tightness and refinements for ISTA on l1-regularized PageRank.
%
% Primary source for the formulation, Algorithm 3, monotone-support theorem,
% support-volume bound, and published work upper bound:
% Fountoulakis et al., Mathematical Programming 174:553--573, 2019.
% Verified against arXiv:1602.01886, especially equations (1), (8)--(10),
% Algorithms 3--4, and Theorems 1--3 (published-paper Sections 4--5).
%
% Scope: the center-star and direct-sum lower bounds, the residual-contraction
% upper bounds, and the thresholded coordinate method below are new arguments.
% They are deliberately separated from statements quoted from the source.
% The relative stopping parameter is denoted by delta, rather than epsilon or
% eta, to avoid collision with the manuscript's PageRank tolerances and the
% FISTA/Catalyst parameters.

\section{Tightness and Refinements for Local ISTA on
\texorpdfstring{$\ell_1$}{l1}-Regularized PageRank}
\label{sec:rppr-ista-tightness}

This section determines which parts of the classical local-ISTA work bound
for \(\ell_1\)-regularized PageRank are genuine.  There are three distinct
conclusions.  First, the factor \(1/(\alpha\rho)\) and the support-volume
bound \(1/\rho\) are simultaneously attained by a center-seeded star.
Second, the factor \(\log(1/\alpha)\) in the published ISTA analysis is not
necessary: a direct contraction argument removes \(\alpha\) from inside the
logarithm.  Third, a remaining \(\log(1/\rho)\) factor is necessary for the
original full-batch ISTA in the general teleportation-distribution setting,
but it is not intrinsic to the variational problem.  A residual-thresholded
coordinate implementation removes that logarithm for constant relative
accuracy.

The distinction between an algorithm-specific lower bound and a
problem-level lower bound is essential.  Nothing in this section proves an
\(\Omega(1/(\alpha\rho))\) lower bound for every local solver of the
regularized problem.

\subsection{ISTA, the relative KKT residual, and the work model}
\label{subsec:rppr-ista-setup}

We use the source-aligned formulation of
\cref{subsec:shared-source-aligned-problem}, in particular
\cref{eq:shared-seed,eq:shared-pagerank-matrices,eq:shared-rppr-objective}.
Throughout this section we retain the source theorem's additional restrictions
\(\alpha\in(0,1)\) and \(0<\rho\leq\norm{s}_\infty\).  Define
\begin{equation}
    a:=\frac{1-\alpha}{2},
    \qquad
    c:=\frac{1+\alpha}{2}=1-a,
    \qquad
    B:=D^{-1/2}AD^{-1/2},
    \qquad
    P:=D^{-1}A.
    \label{eq:rppr-ista-matrices}
\end{equation}
Then \(Q=cI-aB\), \(P\) is row stochastic, and
\begin{equation}
    QD^{1/2}\one=\alpha D^{1/2}\one.
    \label{eq:rppr-degree-eigenvector}
\end{equation}

Using the shared regularizer \(g_\rho\), let
\begin{equation}
    \mathcal{T}_\rho(x)
    :=
    \prox_{g_\rho}\bigl(x-\nabla f(x)\bigr)
    =
    \prox_{g_\rho}\bigl((I-Q)x+\alpha D^{-1/2}s\bigr).
    \label{eq:rppr-ista-map}
\end{equation}
The full-batch ISTA iterates are
\begin{equation}
    x_0=0,
    \qquad
    x_{k+1}=\mathcal{T}_\rho(x_k).
    \label{eq:rppr-batch-ista}
\end{equation}
Under the initialization and nonnegative seed above,
\citet[Theorem~1]{fountoulakis2019variational} proves
\begin{equation}
    0\leq x_k\leq x_{k+1}\leq x^\star,
    \qquad
    \supp(x_k)\subseteq\supp(x_{k+1})\subseteq S^\star,
    \qquad
    \nabla f(x_k)\leq0,
    \label{eq:rppr-monotone-invariants}
\end{equation}
where \(x^\star:=x^\star(\rho)\) and
\(S^\star:=S^\star(\rho)\) use the shared RPPR notation.  The same
source proves
\begin{equation}
    \vol(S^\star)\leq\frac1\rho.
    \label{eq:rppr-support-upper}
\end{equation}

The source writes its relative stopping parameter with an undifferentiated
epsilon glyph. We rename
it \(\delta\in(0,1)\) and define the degree-normalized negative gradient,
the KKT excess, and its maximum norm by
\begin{equation}
    r(x):=-D^{-1/2}\nabla f(x),
    \qquad
    h(x):=\bigl(r(x)-\alpha\rho\one\bigr)_+,
    \qquad
    \mathcal{R}_\rho(x):=\norm{h(x)}_\infty.
    \label{eq:rppr-kkt-excess}
\end{equation}
Here and below the positive part is coordinatewise.  Since
\(\nabla f(x_k)\leq0\), Algorithm~3 of
\citet{fountoulakis2019variational} stops precisely when
\begin{equation}
    \norm{D^{-1/2}\nabla f(x_k)}_\infty
    \leq(1+\delta)\alpha\rho
    \quad\Longleftrightarrow\quad
    \mathcal{R}_\rho(x_k)\leq\delta\alpha\rho.
    \label{eq:rppr-relative-stopping}
\end{equation}
Thus \(\delta\) is dimensionless; the absolute KKT-excess tolerance is
\(\delta\alpha\rho\).

For the batch method, we use the degree-weighted work measure already defined
in \cref{eq:source-degree-work}:
\begin{equation}
    W_{\mathrm{batch}}(T)
    :=
    \sum_{k=0}^{T-1}
    \left(
        \vol(\supp(x_k))+\vol(\supp(x_{k+1}))
    \right).
    \label{eq:rppr-batch-work}
\end{equation}
On an unweighted graph this is equivalent, up to a universal constant, to the
\(|S^\star|+\widehat{\vol}(S^\star)\) per-iteration accounting in
\citet[Theorem~3]{fountoulakis2019variational}.  In particular,
\cref{eq:rppr-monotone-invariants,eq:rppr-support-upper} imply
\begin{equation}
    W_{\mathrm{batch}}(T)\leq\frac{2T}{\rho}.
    \label{eq:rppr-work-from-iterations}
\end{equation}

\begin{lemma}[Increment--residual identity]
\label{lem:rppr-increment-residual}
For every batch-ISTA iterate,
\begin{equation}
    x_{k+1}-x_k=D^{1/2}h(x_k).
    \label{eq:rppr-increment-residual}
\end{equation}
Consequently, \cref{eq:rppr-relative-stopping} is equivalent to
\begin{equation}
    \norm{D^{-1/2}(x_{k+1}-x_k)}_\infty
    \leq\delta\alpha\rho.
    \label{eq:rppr-step-stopping}
\end{equation}
\end{lemma}

\begin{proof}
Weighted soft thresholding and
\cref{eq:rppr-monotone-invariants} give, coordinatewise,
\[
    x_{k+1,i}
    =
    \max\left\{
        x_{k,i}-\nabla_i f(x_k)-\alpha\rho\sqrt{d_i},0
    \right\}.
\]
Every positive coordinate of \(x_k\) satisfies
\(-\nabla_i f(x_k)\geq\alpha\rho\sqrt{d_i}\), while a zero coordinate is
activated exactly when this inequality is strict.  Subtracting \(x_{k,i}\)
therefore yields
\[
    x_{k+1,i}-x_{k,i}
    =
    \sqrt{d_i}
    \left(
        -\frac{\nabla_i f(x_k)}{\sqrt{d_i}}-\alpha\rho
    \right)_+,
\]
which is \cref{eq:rppr-increment-residual}.  The stopping equivalence follows
immediately.
\end{proof}

\subsection{What the stopping rule certifies for unregularized PageRank}
\label{subsec:rppr-ppr-certificate}

Recall from \cref{eq:shared-ppr-solution} that \(x^0\) is the
unregularized minimizer and \(\pi=D^{1/2}x^0\) is its lazy personalized
PageRank vector.  For an approximate regularized solution \(x\), write
\(p:=D^{1/2}x\).

\begin{proposition}[Regularized stopping implies approximate PPR]
\label{prop:rppr-to-ppr-certificate}
Suppose \(\nabla f(x)\leq0\) and
\(\mathcal{R}_\rho(x)\leq\delta\alpha\rho\).  Then
\begin{equation}
    0
    \leq
    D^{-1}(\pi-p)
    =D^{-1/2}(x^0-x)
    \leq
    (1+\delta)\rho\one.
    \label{eq:rppr-to-ppr-coordinate-bound}
\end{equation}
In particular,
\begin{equation}
    \norm{D^{-1}(\pi-p)}_\infty
    \leq(1+\delta)\rho.
    \label{eq:rppr-to-ppr-infinity-bound}
\end{equation}
Thus a target PPR accuracy \(\epsppr\) is guaranteed by
choosing
\begin{equation}
    (1+\delta)\rho\leq\epsppr.
    \label{eq:rppr-rho-delta-choice}
\end{equation}
\end{proposition}

\begin{proof}
Let \(z:=-\nabla f(x)\).  The assumptions imply
\[
    0\leq z\leq(1+\delta)\alpha\rho D^{1/2}\one.
\]
Moreover,
\[
    Q(x^0-x)=z.
\]
The matrix \(Q=cI-aB\) is a nonsingular symmetric \(M\)-matrix.  Indeed,
\(a/c<1\), \(B\geq0\), and
\[
    Q^{-1}
    =
    \frac1c\sum_{t=0}^{\infty}\left(\frac ac B\right)^t
    \geq0.
\]
Using this entrywise nonnegativity together with
\cref{eq:rppr-degree-eigenvector} gives
\[
    0
    \leq x^0-x
    =Q^{-1}z
    \leq
    (1+\delta)\alpha\rho Q^{-1}D^{1/2}\one
    =(1+\delta)\rho D^{1/2}\one.
\]
Multiplication by \(D^{-1/2}\) proves the result.
\end{proof}

\begin{remark}[Approximation, not exact equality]
\label{rem:rppr-not-exact-ppr}
For fixed \(\rho>0\), arbitrarily accurate optimization converges to
\(x^\star(\rho)\), which generally differs from \(x^0\).  The guarantee
in \cref{prop:rppr-to-ppr-certificate} is an approximation statement.  Exact
equality requires \(\rho=0\), a special instance on which the regularizer is
inactive, or a subsequent unregularized debiasing solve.
\end{remark}

\subsection{A sharper upper bound for the same full-batch ISTA}
\label{subsec:rppr-sharp-batch-upper}

The logarithm in \citet[Theorem~3]{fountoulakis2019variational} enters by
first reducing the objective gap to order
\(\delta^2\rho^2\alpha^2d_{\min}\) and then converting objective gap to the
stopping residual.  The next lemma contracts the stopping residual directly.

\begin{lemma}[Direct contraction of the KKT excess]
\label{lem:rppr-kkt-contraction}
Let \(h_k:=h(x_k)\).  Then
\begin{equation}
    0\leq h_{k+1}\leq a(I+P)h_k
    \label{eq:rppr-kkt-vector-contraction}
\end{equation}
coordinatewise, and hence
\begin{equation}
    \norm{h_{k+1}}_\infty
    \leq(1-\alpha)\norm{h_k}_\infty.
    \label{eq:rppr-kkt-infinity-contraction}
\end{equation}
\end{lemma}

\begin{proof}
By \cref{lem:rppr-increment-residual},
\(x_{k+1}-x_k=D^{1/2}h_k\).  Therefore
\begin{align}
    r(x_{k+1})
    &=r(x_k)-D^{-1/2}QD^{1/2}h_k \notag\\
    &=r(x_k)-(cI-aP)h_k.
    \label{eq:rppr-r-update}
\end{align}
The definition of the positive part gives
\(r(x_k)\leq\alpha\rho\one+h_k\).  Subtracting
\(\alpha\rho\one\) from \cref{eq:rppr-r-update} thus yields
\[
    r(x_{k+1})-\alpha\rho\one
    \leq
    (1-c)h_k+aPh_k
    =a(I+P)h_k.
\]
The right-hand side is nonnegative, so taking positive parts proves
\cref{eq:rppr-kkt-vector-contraction}.  Since \(P\) is row stochastic,
\(\norm{I+P}_{\infty}=2\), and \(2a=1-\alpha\), proving
\cref{eq:rppr-kkt-infinity-contraction}.
\end{proof}

Define
\begin{equation}
    H_\rho(s)
    :=
    \norm{\bigl(D^{-1}s-\rho\one\bigr)_+}_\infty,
    \qquad
    \log_+(z):=\max\{\log z,0\}.
    \label{eq:rppr-H-definition}
\end{equation}

\begin{theorem}[Sharpened residual and work upper bounds]
\label{thm:rppr-sharp-batch-upper}
Let \(T\) be the first index satisfying
\cref{eq:rppr-relative-stopping}.  Then
\begin{equation}
    T
    \leq
    \left\lceil
        \frac{
            \log_+\bigl(H_\rho(s)/(\delta\rho)\bigr)
        }{-\log(1-\alpha)}
    \right\rceil
    \leq
    \left\lceil
        \frac1\alpha
        \log_+\left(\frac{H_\rho(s)}{\delta\rho}\right)
    \right\rceil.
    \label{eq:rppr-sharp-iteration-upper}
\end{equation}
Consequently,
\begin{equation}
    W_{\mathrm{batch}}(T)
    \leq
    \frac{2}{\rho}
    \left\lceil
        \frac{
            \log_+\bigl(H_\rho(s)/(\delta\rho)\bigr)
        }{-\log(1-\alpha)}
    \right\rceil.
    \label{eq:rppr-sharp-work-upper}
\end{equation}
In particular, because \(H_\rho(s)\leq1/d_{\min}\),
\begin{equation}
    W_{\mathrm{batch}}(T)
    =
    O\left(
        \frac{1}{\alpha\rho}
        \left[1+\log_+\left(
            \frac{1}{\delta\rho d_{\min}}
        \right)\right]
    \right).
    \label{eq:rppr-sharp-worst-upper}
\end{equation}
The teleportation parameter \(\alpha\) does not occur inside the logarithm.
\end{theorem}

\begin{proof}
At initialization,
\begin{equation}
    h_0
    =
    \alpha\bigl(D^{-1}s-\rho\one\bigr)_+,
    \qquad
    \norm{h_0}_\infty=\alpha H_\rho(s).
    \label{eq:rppr-initial-excess}
\end{equation}
Iterating \cref{eq:rppr-kkt-infinity-contraction} gives
\[
    \mathcal{R}_\rho(x_k)
    =\norm{h_k}_\infty
    \leq
    \alpha(1-\alpha)^kH_\rho(s).
\]
The right-hand side is at most \(\delta\alpha\rho\) once the first bound in
\cref{eq:rppr-sharp-iteration-upper} holds.  The second bound uses
\(-\log(1-\alpha)\geq\alpha\).  The work bound follows from
\cref{eq:rppr-work-from-iterations}.  Finally,
\(s_i/d_i\leq1/d_{\min}\) for every \(i\), which proves
\cref{eq:rppr-sharp-worst-upper}.
\end{proof}

The restricted strong-convexity parameter can also be retained in this
sharper analysis.  Define
\begin{equation}
    \mu_\star:=\lambda_{\min}(Q_{S^\star S^\star}),
    \qquad
    B_\rho(s)
    :=
    \norm{D^{-1/2}(s-\rho d)_+}_2,
    \label{eq:rppr-restricted-parameters}
\end{equation}
where \(d=D\one\).

\begin{proposition}[Restricted contraction without \(\log(1/\alpha)\)]
\label{prop:rppr-restricted-contraction}
Assume \(S^\star\neq\emptyset\).  Then
\begin{equation}
    T
    \leq
    \left\lceil
        \frac{
            \log_+\left(
                B_\rho(s)/(\delta\rho\sqrt{d_{\min}})
            \right)
        }{-\log(1-\mu_\star)}
    \right\rceil,
    \label{eq:rppr-restricted-iteration-upper}
\end{equation}
and therefore
\begin{equation}
    W_{\mathrm{batch}}(T)
    =
    O\left(
        \frac{1}{\rho\mu_\star}
        \left[
            1+
            \log_+\left(
                \frac{B_\rho(s)}
                     {\delta\rho\sqrt{d_{\min}}}
            \right)
        \right]
    \right).
    \label{eq:rppr-restricted-work-upper}
\end{equation}
\end{proposition}

\begin{proof}
Let \(\Delta_k:=x_{k+1}-x_k\).  Soft thresholding is coordinatewise
nonexpansive.  Both \(\mathcal{T}_\rho(x_{k+1})\) and
\(\mathcal{T}_\rho(x_k)\) are supported on \(S^\star\) by
\cref{eq:rppr-monotone-invariants}; hence
\begin{align}
    \norm{\Delta_{k+1}}_2
    &\leq
    \norm{(I-Q_{S^\star S^\star})\Delta_k}_2 \notag\\
    &\leq
    (1-\mu_\star)\norm{\Delta_k}_2.
    \label{eq:rppr-restricted-step-contraction}
\end{align}
Here all eigenvalues of \(Q_{S^\star S^\star}\) lie in
\([\mu_\star,1]\).  Moreover,
\[
    \Delta_0
    =
    \alpha D^{-1/2}(s-\rho d)_+,
    \qquad
    \norm{\Delta_0}_2=\alpha B_\rho(s).
\]
By \cref{eq:rppr-step-stopping}, termination is guaranteed whenever
\[
    \norm{\Delta_k}_2
    \leq\delta\alpha\rho\sqrt{d_{\min}}.
\]
The factor \(\alpha\) cancels between the initial step and the stopping
threshold, giving \cref{eq:rppr-restricted-iteration-upper}.  The work bound
again follows from \cref{eq:rppr-work-from-iterations} and
\(-\log(1-\mu_\star)\geq\mu_\star\).
\end{proof}

\begin{remark}[Comparison with the published bound]
\label{rem:rppr-published-log-comparison}
After renaming the source accuracy parameter as \(\delta\), the published
unweighted-graph work bound has the form
\begin{equation}
    O\left(
        \frac{1}{\rho\mu}
        \log\frac{2}{\delta^2\rho^2\alpha^2d_{\min}}
    \right),
    \label{eq:rppr-published-work-bound}
\end{equation}
where \(\mu\geq\alpha\) is a lower bound on the restricted curvature.
\Cref{thm:rppr-sharp-batch-upper,prop:rppr-restricted-contraction} prove that
the \(\log(1/\alpha)\) contribution is an artifact of the objective-gap to
residual conversion.  They do not, by themselves, remove
\(\log(1/\rho)\) for arbitrary seed distributions.
\end{remark}

\subsection{A center-star lower bound}
\label{subsec:rppr-center-star-lower}

We first show that one center-seeded star simultaneously saturates the
support-volume bound and the polynomial factors in the full-batch ISTA work
bound.

\begin{theorem}[Center-star lower bound for full-batch ISTA]
\label{thm:rppr-center-star-lower}
Fix
\begin{equation}
    0<\alpha\leq\frac14,
    \qquad
    0<\rho\leq\frac1{16},
    \qquad
    0<\delta<1.
    \label{eq:rppr-star-parameter-range}
\end{equation}
Let \(G=K_{1,m}\), put the seed at its center \(v\), and choose
\begin{equation}
    m:=\floor{\frac1{8\rho}},
    \qquad
    s=e_v.
    \label{eq:rppr-star-size}
\end{equation}
Then
\begin{equation}
    S^\star=V,
    \qquad
    \vol(S^\star)=2m=\Theta(1/\rho),
    \label{eq:rppr-star-full-support}
\end{equation}
and full-batch ISTA requires degree-weighted work
\begin{equation}
    W_{\mathrm{batch}}(T)
    \geq
    \frac{3}{32\alpha\rho}\log\frac2\delta.
    \label{eq:rppr-star-work-lower}
\end{equation}
Consequently, on this single-seed instance,
\begin{equation}
    W_{\mathrm{batch}}(T)
    =
    \Theta\left(
        \frac{1}{\alpha\rho}
        \left[1+\log\frac1\delta\right]
    \right).
    \label{eq:rppr-star-tight-order}
\end{equation}
\end{theorem}

\begin{proof}
Let \(\ell\) be any leaf and define
\begin{equation}
    \kappa:=\rho m,
    \qquad
    u_k:=\sqrt{m}\,x_{k,v},
    \qquad
    v_k:=m x_{k,\ell}.
    \label{eq:rppr-star-aggregate-coordinates}
\end{equation}
The choice of \(m\) and \(\rho\leq1/16\) gives
\begin{equation}
    \frac1{16}\leq\kappa\leq\frac18,
    \qquad
    2m\geq\frac1{8\rho}.
    \label{eq:rppr-kappa-range}
\end{equation}
Symmetry is preserved by ISTA.  Applying
\cref{eq:rppr-ista-map} at the center and at a leaf gives
\begin{align}
    u_{k+1}
    &=a(u_k+v_k)+\alpha(1-\kappa),
    \label{eq:rppr-star-u-recurrence}\\
    v_{k+1}
    &=\bigl[a(u_k+v_k)-\alpha\kappa\bigr]_+.
    \label{eq:rppr-star-v-recurrence}
\end{align}
In particular,
\[
    u_1=\alpha(1-\kappa),
    \qquad
    v_1=0.
\]
All leaves activate in the next step, because
\begin{equation}
    a(1-\kappa)
    \geq\frac38\cdot\frac78
    =\frac{21}{64}
    >\frac18
    \geq\kappa.
    \label{eq:rppr-star-leaf-activation}
\end{equation}
Thereafter the positive part in
\cref{eq:rppr-star-v-recurrence} is inactive.  With
\(t_k:=u_k+v_k\), we obtain, for every \(k\geq1\),
\begin{equation}
    t_{k+1}
    =(1-\alpha)t_k+\alpha(1-2\kappa).
    \label{eq:rppr-star-t-recurrence}
\end{equation}
The fixed point satisfies
\begin{equation}
    t^\star=1-2\kappa,
    \qquad
    u^\star-v^\star=\alpha,
    \label{eq:rppr-star-fixed-relations}
\end{equation}
and therefore
\begin{equation}
    u^\star=\frac{1-2\kappa+\alpha}{2},
    \qquad
    v^\star=\frac{1-2\kappa-\alpha}{2}>0.
    \label{eq:rppr-star-fixed-point}
\end{equation}
The last inequality follows from
\(2\kappa+\alpha\leq1/2\).  Thus every coordinate of the unique fixed point
of ISTA is positive, proving \(S^\star=V\).  Since the star has volume
\(2m\), \cref{eq:rppr-star-full-support} follows.

Let
\begin{equation}
    G_0
    :=t^\star-t_1
    =1-2\kappa-\alpha(1-\kappa).
    \label{eq:rppr-star-G0}
\end{equation}
The parameter range gives \(G_0\geq1/2\), and
\cref{eq:rppr-star-t-recurrence} gives the exact slow mode
\begin{equation}
    t^\star-t_k
    =G_0(1-\alpha)^{k-1},
    \qquad k\geq1.
    \label{eq:rppr-star-slow-mode}
\end{equation}
Moreover, \(u_k-v_k=\alpha\) for every \(k\geq2\).  Hence, for
\(k\geq2\),
\begin{equation}
    u_{k+1}-u_k
    =v_{k+1}-v_k
    =\frac{\alpha G_0}{2}(1-\alpha)^{k-1}.
    \label{eq:rppr-star-equal-increments}
\end{equation}
Using \cref{lem:rppr-increment-residual} and the definitions of \(u_k,v_k\),
we obtain the same normalized KKT excess at the center and at every leaf:
\begin{equation}
    \mathcal{R}_\rho(x_k)
    =
    \frac{\alpha G_0}{2m}(1-\alpha)^{k-1},
    \qquad k\geq2.
    \label{eq:rppr-star-exact-residual}
\end{equation}

Let \(T\) be the first stopping index.  The activation inequality
\cref{eq:rppr-star-leaf-activation} is quantitatively stronger than mere
activation: at \(x_1\), the excess at every leaf is
\[
    h_\ell(x_1)
    =
    \alpha\rho\,
    \frac{a(1-\kappa)-\kappa}{\kappa}
    \geq
    \frac{13}{8}\alpha\rho
    >
    \delta\alpha\rho.
\]
Thus \(T\geq2\).  From
\cref{eq:rppr-relative-stopping,eq:rppr-star-exact-residual},
\begin{equation}
    G_0(1-\alpha)^{T-1}
    \leq2\kappa\delta.
    \label{eq:rppr-star-stop-inequality}
\end{equation}
Therefore
\begin{equation}
    T-1
    \geq
    \frac{\log(G_0/(2\kappa\delta))}
         {-\log(1-\alpha)}
    \geq
    \frac{3}{4\alpha}\log\frac2\delta,
    \label{eq:rppr-star-iteration-lower}
\end{equation}
where we used \(G_0\geq1/2\), \(2\kappa\leq1/4\), and
\(-\log(1-\alpha)\leq\alpha/(1-\alpha)\leq4\alpha/3\).

For every update \(k=1,\ldots,T-1\), the new iterate \(x_{k+1}\) has full
support and hence contributes \(\vol(V)=2m\) to
\cref{eq:rppr-batch-work}.  Consequently,
\[
    W_{\mathrm{batch}}(T)
    \geq2m(T-1)
    \geq
    \frac1{8\rho}\cdot
    \frac{3}{4\alpha}\log\frac2\delta,
\]
which is \cref{eq:rppr-star-work-lower}.  On this instance
\(H_\rho(s)=1/m-\rho=\Theta(\rho)\), so
\cref{thm:rppr-sharp-batch-upper} supplies the matching upper order in
\cref{eq:rppr-star-tight-order}.
\end{proof}

\subsection{The remaining logarithm for general full-batch ISTA}
\label{subsec:rppr-direct-sum-lower}

The center star cannot force \(\log(1/\rho)\), because its initial
degree-normalized seed density is already \(1/m=\Theta(\rho)\).  For the
general seed-distribution setting, a direct sum separates the large active
volume from a constant-scale slow residual.

\begin{theorem}[Direct-sum lower bound]
\label{thm:rppr-direct-sum-lower}
Under \cref{eq:rppr-star-parameter-range}, let
\begin{equation}
    G=K_2\mathbin{\dot\cup}K_{1,m},
    \qquad
    m=\floor{\frac1{8\rho}},
    \qquad
    s=\frac12e_u+\frac12e_w,
    \label{eq:rppr-direct-sum-instance}
\end{equation}
where \(u\) is one endpoint of \(K_2\) and \(w\) is the star center.  Then
\(S^\star=V\), \(\vol(S^\star)=\Theta(1/\rho)\), and
\begin{equation}
    W_{\mathrm{batch}}(T)
    \geq
    \frac{3}{32\alpha\rho}
    \log\left(\frac{1}{8\delta\rho}\right).
    \label{eq:rppr-direct-sum-work-lower}
\end{equation}
Consequently, if arbitrary teleportation distributions and disconnected
graphs are permitted, the worst-case work of the original full-batch ISTA is
\begin{equation}
    W_{\mathrm{batch}}^{\mathrm{worst}}(\alpha,\rho,\delta)
    =
    \Theta\left(
        \frac{1}{\alpha\rho}
        \left[1+\log\frac{1}{\delta\rho}\right]
    \right)
    \label{eq:rppr-batch-general-tight}
\end{equation}
throughout this parameter regime.
\end{theorem}

\begin{proof}
Each component has seed mass \(1/2\).  On the star, the aggregate recurrence
from
\cref{eq:rppr-star-u-recurrence,eq:rppr-star-v-recurrence} becomes
\begin{align}
    u_{k+1}&=a(u_k+v_k)+\alpha(1/2-\kappa),
    \label{eq:rppr-half-star-u}\\
    v_{k+1}&=\bigl[a(u_k+v_k)-\alpha\kappa\bigr]_+.
    \label{eq:rppr-half-star-v}
\end{align}
All leaves activate in the second iterate because
\begin{equation}
    a(1/2-\kappa)
    \geq\frac38\left(\frac12-\frac18\right)
    =\frac9{64}
    >\frac18
    \geq\kappa.
    \label{eq:rppr-half-star-activation}
\end{equation}
The star fixed point satisfies
\[
    u^\star+v^\star=1/2-2\kappa,
    \qquad
    u^\star-v^\star=\alpha/2,
\]
so
\[
    v^\star
    =\frac{(1-\alpha)/2-2\kappa}{2}
    \geq\frac{1/8}{2}>0.
\]
Thus the entire star belongs to \(S^\star\).

Let \(x_k^{u},x_k^{v}\) denote the two degree-one coordinates on the
\(K_2\) component.  We have
\[
    x_1^{u}=\alpha(1/2-\rho),
    \qquad
    x_1^{v}=0.
\]
The unseeded endpoint activates in the next iterate because
\[
    a(1/2-\rho)
    \geq\frac38\left(\frac12-\frac1{16}\right)
    >\frac1{16}
    \geq\rho.
\]
Thereafter, with \(z_k:=x_k^{u}+x_k^{v}\),
\begin{equation}
    z_{k+1}
    =(1-\alpha)z_k+\alpha(1/2-2\rho),
    \qquad k\geq1.
    \label{eq:rppr-K2-sum-recurrence}
\end{equation}
The fixed point obeys
\[
    z^\star=1/2-2\rho,
    \qquad
    x^{u,\star}-x^{v,\star}=\alpha/2,
\]
and hence
\[
    x^{v,\star}
    =\frac{(1-\alpha)/2-2\rho}{2}>0.
\]
Thus \(S^\star=V\), and its volume is \(2m+2=\Theta(1/\rho)\).

Define \(E_k:=z^\star-z_k\).  From
\cref{eq:rppr-K2-sum-recurrence},
\begin{equation}
    E_k=E_1(1-\alpha)^{k-1},
    \qquad
    E_1=1/2-2\rho-\alpha(1/2-\rho)
    \geq\frac{17}{64}>\frac14.
    \label{eq:rppr-K2-error}
\end{equation}
For every \(k\geq2\), the difference
\(x_k^{u}-x_k^{v}=\alpha/2\) is fixed, so the two increments are
equal.
By \cref{lem:rppr-increment-residual}, their KKT excess is
\begin{equation}
    h_u(x_k)=h_v(x_k)
    =\frac{\alpha E_k}{2}.
    \label{eq:rppr-K2-excess}
\end{equation}
Global termination therefore requires \(E_T\leq2\delta\rho\), and hence
\begin{equation}
    T-1
    \geq
    \frac{\log(E_1/(2\delta\rho))}{-\log(1-\alpha)}
    \geq
    \frac{3}{4\alpha}
    \log\left(\frac{1}{8\delta\rho}\right).
    \label{eq:rppr-direct-sum-iteration-lower}
\end{equation}
From the second iterate onward, the star has full support.  Exactly as in the
last paragraph of the proof of \cref{thm:rppr-center-star-lower},
\[
    W_{\mathrm{batch}}(T)
    \geq2m(T-1),
\]
which proves \cref{eq:rppr-direct-sum-work-lower}.  The matching worst-case
upper bound in \cref{eq:rppr-batch-general-tight} follows from
\cref{thm:rppr-sharp-batch-upper}, because
\(H_\rho(s)\leq1\) and \(\vol(S^\star)\leq1/\rho\) on every unweighted
instance.
\end{proof}

\begin{remark}[What the direct sum does not prove]
\label{rem:rppr-connected-single-seed-open}
\Cref{thm:rppr-direct-sum-lower} uses two seeded components.  It proves that
\(\log(1/\rho)\) is necessary for the general full-batch theorem, but it does
not prove the same logarithm for connected graphs with a single-node seed.
The center star already proves the tight polynomial factor
\(1/(\alpha\rho)\) in that narrower setting.  Whether a connected,
single-seed construction forces the additional \(\log(1/\rho)\) for
full-batch ISTA remains open here.
\end{remark}

\subsection{A logarithm-free local coordinate implementation}
\label{subsec:rppr-threshold-coordinate}

The logarithm in \cref{thm:rppr-direct-sum-lower} comes from repeatedly
scanning the already-active star while the small \(K_2\) component is still
converging.  A coordinate method need not perform those stale scans.

\begin{algorithm}[t]
\caption{Residual-thresholded coordinate ISTA}
\label{alg:rppr-threshold-coordinate}
\begin{algorithmic}[1]
\REQUIRE \(G,s,\alpha,\rho,\delta\), with \(\delta\in(0,1]\)
\STATE \(x\gets0\), \(\nabla f(x)\gets-\alpha D^{-1/2}s\)
\STATE Initialize a local queue with coordinates in \(\supp(s)\) satisfying
       \(h_i(x)\geq\delta\alpha\rho\)
\WHILE{the queue is nonempty}
    \STATE Remove any queued coordinate \(i\) that still satisfies
           \(h_i(x)\geq\delta\alpha\rho\))
    \STATE \(\Delta\gets\sqrt{d_i}\,h_i(x)\)
    \STATE \(x_i\gets x_i+\Delta\)
    \STATE Update \(\nabla f(x)\gets\nabla f(x)+\Delta Qe_i\) on
           \(\{i\}\cup\mathcal{N}(i)\))
    \STATE Refresh the queue status of \(i\) and its neighbors
\ENDWHILE
\RETURN \(x\)
\end{algorithmic}
\end{algorithm}

Only the selected adjacency list is scanned.  Since the initial negative
gradient is supported on \(\supp(s)\) and a coordinate update changes the
gradient only at that coordinate and its neighbors, the violating set can be
maintained without a global graph scan.  We assume that \(s\) is supplied in
sparse form; otherwise, an \(O(\nnz(s))\) input-initialization cost must be
added to the work below.

\begin{theorem}[Coordinate work bound]
\label{thm:rppr-threshold-coordinate-work}
Algorithm~\ref{alg:rppr-threshold-coordinate} terminates with
\begin{equation}
    \nabla f(x)\leq0,
    \qquad
    \mathcal{R}_\rho(x)\leq\delta\alpha\rho,
    \qquad
    0\leq x\leq x^\star.
    \label{eq:rppr-coordinate-certificate}
\end{equation}
Its total adjacency-list work satisfies
\begin{equation}
    W_{\mathrm{coord}}
    :=\sum_{\text{updates at }i}d_i
    \leq\frac{1}{\delta\alpha\rho}.
    \label{eq:rppr-coordinate-work-bound}
\end{equation}
For constant \(\delta\), this is
\(O(1/(\alpha\rho))\), with no logarithmic factor.
\end{theorem}

\begin{proof}
First, the invariants in \cref{eq:rppr-coordinate-certificate} are preserved.
Suppose coordinate \(i\) is selected.  Before the update,
\[
    \nabla_i f(x)
    =-\sqrt{d_i}\bigl(\alpha\rho+h_i(x)\bigr).
\]
Since \(Q_{ii}=c\) and all off-diagonal entries of \(Q\) are nonpositive,
\begin{align}
    \nabla_i f(x+\Delta e_i)
    &\leq
    -\sqrt{d_i}\bigl(\alpha\rho+h_i(x)\bigr)
    +c\sqrt{d_i}h_i(x) \notag\\
    &=-\sqrt{d_i}
       \bigl(\alpha\rho+(1-c)h_i(x)\bigr)<0.
    \label{eq:rppr-coordinate-gradient-invariant}
\end{align}
The gradients at neighboring coordinates only decrease, and all remaining
gradients are unchanged.  Hence \(\nabla f(x)\leq0\) is invariant.

Next, \(I-Q=a(I+B)\) is entrywise nonnegative, so the map
\(\mathcal{T}_\rho\) in \cref{eq:rppr-ista-map} is monotone in the
coordinatewise order.  If \(0\leq x\leq x^\star\), then
\[
    \mathcal{T}_\rho(x)
    \leq\mathcal{T}_\rho(x^\star)=x^\star.
\]
The coordinate update sets
\(x_i\) to \([\mathcal{T}_\rho(x)]_i\), while leaving all other coordinates
unchanged; the selected violation makes this new value at least \(x_i\).
Induction from \(x=0\) proves \(0\leq x\leq x^\star\), and in particular the
method never activates a coordinate outside \(S^\star\).

It remains to bound work.  Define the nonnegative residual-mass potential
\begin{equation}
    \Phi(x)
    :=\sum_{j\in V}d_jr_j(x)
    =-\left\langle D^{1/2}\one,\nabla f(x)\right\rangle.
    \label{eq:rppr-coordinate-potential}
\end{equation}
The initialization gives \(\Phi(0)=\alpha\).  An update at \(i\) changes the
potential by the exact amount
\begin{align}
    \Phi(x+\Delta e_i)-\Phi(x)
    &=-\Delta\left\langle D^{1/2}\one,Qe_i\right\rangle \notag\\
    &=-\alpha d_i h_i(x),
    \label{eq:rppr-coordinate-potential-drop}
\end{align}
where \cref{eq:rppr-degree-eigenvector} and
\(\Delta=\sqrt{d_i}h_i(x)\) were used.  Every selected coordinate satisfies
\(h_i(x)\geq\delta\alpha\rho\), and hence decreases \(\Phi\) by at least
\(\delta\alpha^2\rho d_i\).  Since \(\Phi\geq0\), summing these decreases
gives
\[
    \delta\alpha^2\rho
    \sum_{\text{updates at }i}d_i
    \leq\Phi(0)=\alpha,
\]
which is \cref{eq:rppr-coordinate-work-bound}.  At termination no coordinate
violates the threshold, giving the remaining statement in
\cref{eq:rppr-coordinate-certificate}.
\end{proof}

For very small \(\delta\), the \(1/\delta\) dependence in
\cref{eq:rppr-coordinate-work-bound} is avoidable by using coordinate updates
only as a coarse localization phase.

\begin{corollary}[Coarse-to-fine logarithm-free-in-\(\rho\) bound]
\label{cor:rppr-coarse-to-fine}
Run Algorithm~\ref{alg:rppr-threshold-coordinate} first with relative
threshold \(1\), obtaining \(\bar{x}\) with
\(\mathcal{R}_\rho(\bar{x})\leq\alpha\rho\).  Starting from \(\bar{x}\), run
full-batch ISTA until
\(\mathcal{R}_\rho(x)\leq\delta\alpha\rho\).  Then
\begin{equation}
    W_{\mathrm{coarse\mbox{-}to\mbox{-}fine}
    }
    \leq
    \frac{1}{\alpha\rho}
    +
    \frac{2}{\rho}
    \left\lceil
        \frac{\log(1/\delta)}{-\log(1-\alpha)}
    \right\rceil
    =
    O\left(
        \frac{1+\log(1/\delta)}{\alpha\rho}
    \right).
    \label{eq:rppr-coarse-to-fine-work}
\end{equation}
In particular, for constant \(\delta\) the work is
\(O(1/(\alpha\rho))\).
\end{corollary}

\begin{proof}
The first phase costs at most \(1/(\alpha\rho)\) by
\cref{thm:rppr-threshold-coordinate-work}.  Its proof also gives
\(0\leq\bar{x}\leq x^\star\) and \(\nabla f(\bar{x})\leq0\), so the monotone
batch-ISTA invariants continue to hold from this warm start.  Applying
\cref{lem:rppr-kkt-contraction} repeatedly gives
\[
    \mathcal{R}_\rho(x_k)
    \leq(1-\alpha)^k\alpha\rho.
\]
The displayed number of cleanup iterations is therefore sufficient.  Each
cleanup iteration has work at most \(2\vol(S^\star)\leq2/\rho\), proving
\cref{eq:rppr-coarse-to-fine-work}.
\end{proof}

\subsection{Tightness status and remaining boundary}
\label{subsec:rppr-tightness-status}

The results above support the following precise conclusions.

\begin{enumerate}[leftmargin=*]
    \item The solution-volume guarantee
    \(\vol(S^\star)\leq1/\rho\) is tight up to a universal constant, even for
    an unweighted center-seeded star.

    \item For the original full-batch ISTA and general teleportation
    distributions, the worst-case degree-weighted work is
    \[
        \Theta\left(
            \frac{1}{\alpha\rho}
            \left[1+\log\frac{1}{\delta\rho}\right]
        \right).
    \]
    The \(\log(1/\alpha)\) in the published upper bound is unnecessary; the
    \(\log(1/\rho)\) is necessary in this general batch setting.

    \item For a center-seeded star, the exact order is instead
    \[
        \Theta\left(
            \frac{1}{\alpha\rho}
            \left[1+\log\frac1\delta\right]
        \right),
    \]
    because the initial degree-normalized seed density is already
    \(\Theta(\rho)\).

    \item The logarithm in \(1/\rho\) is not intrinsic to solving the
    regularized problem to the relative certificate
    \cref{eq:rppr-relative-stopping}.  Thresholded coordinate ISTA costs
    \(O(1/(\delta\alpha\rho))\), and the coarse-to-fine method costs
    \(O((1+\log(1/\delta))/(\alpha\rho))\).

    \item If \(\delta\) is a fixed constant and the desired unregularized PPR
    accuracy is \(\epsppr\), choosing
    \(\rho=\Theta(\epsppr)\) in
    \cref{eq:rppr-rho-delta-choice} gives
    \(O(1/(\alpha\epsppr))\) coordinate work.

    \item These are algorithm-specific results.  Under an explicit-vector
    output model, the star gives the problem-level output-size obstruction
    \(\Omega(1/\rho)\).  No \(\Omega(1/(\alpha\rho))\) lower bound for every
    local first-order method is proved here.  Indeed, the possibility of
    accelerated sparse methods with improved \(\alpha\)-dependence motivates
    the open problem of \citet{fountoulakis2022open} and the accelerated
    constructions of \citet{martinezrubio2023accelerated}.  Establishing the
    correct local-oracle lower bound remains separate from the batch-ISTA
    tightness theorem above.
\end{enumerate}
